\documentclass[11pt]{article}
\usepackage[margin=1in]{geometry}
\usepackage[T1]{fontenc}
\usepackage{lmodern}
\usepackage{array}
\usepackage{booktabs}
\usepackage{enumitem}

\setlist[itemize]{topsep=3pt,itemsep=2pt,leftmargin=18pt}
\setlength{\parskip}{0.45em}
\setlength{\parindent}{0pt}

\title{ECE51216 Intermediate Report\\Python SAT Solver Project}
\author{Soomin Lee \\ Dogyu Ryu}
\date{Spring 2026}

\begin{document}
\maketitle

\section{Project Overview}

Our project follows Option 1 from the proposal: building a SAT solver based on the DPLL
(Davis--Putnam--Logemann--Loveland) framework for CNF formulas in DIMACS format. The
original goal was to start from a correct baseline DPLL solver and then extend it with more
efficient data structures and heuristics, especially watched literals and VSIDS. The intermediate
goal for this phase was therefore to establish a working baseline implementation that can parse
benchmark instances, perform recursive search, and produce measurable statistics for later
comparison.

\section{Progress Since the Proposal}

The main progress so far has been the implementation of a fresh Python baseline solver rather than
continuing directly from older reference code. This decision was made to keep the code structure
simple, readable, and aligned with our own understanding of the algorithm. The current solver has
already reached the stage where it can read benchmark files, perform search, and report basic
execution statistics.

The implemented components are listed below.

\begin{itemize}
    \item \textbf{DIMACS parser.} We implemented input parsing for CNF problems in DIMACS
    format. The parser handles comment lines, problem headers, standard clause termination with
    \texttt{0}, and SATLIB-style end markers.
    \item \textbf{Baseline DPLL recursion.} The solver performs recursive branching on
    unassigned literals and backtracks when a contradiction is found.
    \item \textbf{Boolean Constraint Propagation (BCP).} Before each branching step, the solver
    repeatedly detects unit clauses and propagates forced assignments.
    \item \textbf{Formula simplification.} After each assignment, clauses satisfied by that literal
    are removed, while falsified literals are deleted from the remaining clauses.
    \item \textbf{Instrumentation.} The implementation records the number of decisions,
    propagations, conflicts, backtracks, recursive calls, maximum recursion depth, and runtime.
\end{itemize}

At this stage, the emphasis has been on correctness and modularity rather than optimization. This
is intentional, because the baseline version will serve as the control point when we later compare
the effect of watched literals and VSIDS.

\section{Current Implementation Status}

The current solver is able to classify both satisfiable and unsatisfiable instances correctly on the
sample benchmark inputs we tested. In addition, it can optionally print a satisfying assignment in a
DIMACS-style literal format when a SAT instance is solved.

To validate the baseline implementation, we ran the solver on three representative examples:

\begin{center}
\begin{tabular}{p{3.1cm} p{1.2cm} p{7.2cm}}
\toprule
\textbf{Instance} & \textbf{Result} & \textbf{Observed Statistics} \\
\midrule
\texttt{example-1.cnf} & SAT & 4 decisions, 13 propagations, 0 conflicts, 0 backtracks, runtime 0.000220 s \\
\texttt{unsat.cnf} & UNSAT & 2 decisions, 2 propagations, 2 conflicts, 2 backtracks, runtime 0.000010 s \\
\texttt{uf20-04.cnf} & SAT & 22 decisions, 67 propagations, 9 conflicts, 16 backtracks, runtime 0.000552 s \\
\bottomrule
\end{tabular}
\end{center}

These results are still preliminary and were obtained only to validate the baseline pipeline from
input parsing to recursive search. They do not yet constitute a full benchmark study, but they do
show that the solver is already functional and ready for the next stage of development.

\section{Changes From the Original Proposal}

The core project objective has not changed: we are still building a SAT solver using the DPLL
framework and planning to evaluate the effect of algorithmic heuristics on performance. However,
two practical changes have been made since the proposal.

First, we decided to use \textbf{Python instead of C/C++} for the implementation. The main reason
is faster prototyping and easier debugging during the development phase. Since the course project
focuses more on algorithmic understanding and comparison than on absolute industrial-level
performance, Python is sufficient for implementing and evaluating the required ideas.

Second, the implementation order has been slightly adjusted. Instead of trying to introduce several
advanced heuristics at once, we first focused on producing a correct and instrumented baseline.
This staged approach makes the later evaluation cleaner, because we will be able to compare:

\begin{itemize}
    \item baseline DPLL,
    \item DPLL with watched literals,
    \item DPLL with watched literals and VSIDS.
\end{itemize}

This comparison structure is consistent with the original evaluation plan, but the development has
been reorganized to prioritize a reliable baseline first.

\section{Remaining Work and Plan}

The next implementation step is to replace naive clause scanning with a watched-literal data
structure so that propagation becomes more efficient. After that, we will add the VSIDS decision
heuristic to improve branching quality. Once these features are integrated, we will run a more
systematic benchmark evaluation using CNF instances from SATLIB and other DIMACS-style test
sets.

Our remaining work is summarized as follows.

\begin{itemize}
    \item Implement watched literals and verify that propagation remains correct.
    \item Add the VSIDS heuristic for decision variable selection.
    \item Expand the benchmark set to include both SAT and UNSAT cases of varying sizes.
    \item Compare solver versions using runtime, number of decisions, number of propagations,
    conflicts, and backtracks.
    \item Summarize the results and discuss which heuristic contributes the most to performance.
\end{itemize}

\section{Conclusion}

In summary, the project has progressed from planning to a working Python baseline SAT solver.
The current implementation already supports DIMACS parsing, DPLL recursion, unit propagation,
and execution statistics, and it has been validated on several sample SAT and UNSAT instances.
Although watched literals and VSIDS are not yet implemented, the project is now at a stable point
where these extensions can be added and evaluated in a structured way.

\end{document}
